\section{Introduction}
\label{sec:intro}

When a deployed model is complemented by human review, the standard question is whether the human knows something the model does not.
A physician examining a chest radiograph may detect subtle findings missed by a convolutional network; a clinician triaging emergency patients may integrate history unavailable to the risk score.
Statistical audits can test for such \emph{residual expertise}---predictive information in the human signal beyond the model \citep{AlurEtAl2023,AlurRaghavanShah2024}.
But for deployment, the sharper question is whether that information changes the decision.

Residual predictive information is not residual decision value.
In a finite-action decision problem governed by a reward matrix, what matters is whether the human's signal moves the posterior across a \emph{decision boundary}.
A physician who shifts the probability of disease from $0.55$ to $0.70$ improves predictions, but if the decision threshold is $0.50$ (the patient is already flagged for admission) or $0.83$ (neither estimate clears the high-confidence bar), the action is unchanged.
The same human signal can carry decision value under one reward structure and none under another---not because the physician's expertise varies, but because the reward geometry determines which posterior movements are consequential.

We formalize this observation through \emph{boundary regret}: the regret of keeping the model's action after observing the human signal (defined formally in Section~\ref{sec:methodology}).
For model belief $b_x \in \Delta^{K-1}$, model-optimal action $a_x = \argmax_a r_a \cdot b_x$, and human-updated belief $b_{x,h} = P(Y \mid x, h)$:
\begin{equation}\label{eq:BR-intro}
  \BR(x, h) \;=\; V(b_{x,h}) \;-\; r_{a_x} \cdot b_{x,h},
\end{equation}
where $V(b) = \max_a r_a \cdot b$ is the terminal value function.
Boundary regret is zero when the model's action remains optimal under the human-updated posterior, and positive when a rival action becomes preferred.

The formal basis is immediate but consequential: in any finite-action Bayesian decision problem, the value of observing the human signal equals expected boundary regret (Theorem~\ref{thm:voi-br}):
\[
  \VoI(H \mid b_x) \;=\; \E_h\bigl[\BR(x, h)\bigr].
\]
The proof is a single application of the Bayesian martingale property; the contribution is not the proof technique but the choice of estimand and its consequences for human--AI audit design.
A companion result (Theorem~\ref{thm:facet-bridge}) provides the geometric characterization: the value of review is strictly positive if and only if some reachable posterior crosses a reward-induced decision boundary.
Signal that improves predictions without crossing any boundary has zero decision value---a formal consequence, not an empirical pattern.
Both results are verified in Lean~4 (Appendix~\ref{app:lean}).

Empirically, we validate this distinction on real and synthetic data across three experiments:
\begin{itemize}[nosep,leftmargin=1.5em]
  \item \textbf{Experiment~1} (synthetic validation): the boundary-regret estimator recovers known values with $\rho = 0.77$ average rank correlation in the boundary-crossing condition and zero false positives across all zero-BR conditions, even when a strong human signal produces large log-loss gains without crossing any decision boundary.
  \item \textbf{Experiment~2} (CheXpert): among the top 10\% of instances by log-loss gain, \textbf{97\%} have zero boundary regret under a false-positive-penalty reward. The same radiologist signal has positive decision value under symmetric accuracy and zero value under asymmetric costs.
  \item \textbf{Experiment~3} (synthetic allocation): on synthetic binary decision problems, an ex ante boundary-regret score allocates scarce review more efficiently than entropy, margin, residual-expertise, and learning-to-defer baselines under asymmetric rewards, while matching them under symmetric accuracy.
\end{itemize}

\noindent\textbf{Contributions.}
\begin{enumerate}[nosep,leftmargin=1.5em]
  \item Boundary regret as a deployment estimand for human--AI complementarity, with formal characterization (Theorems~\ref{thm:voi-br}--\ref{thm:facet-bridge}) and Lean~4 verification.
  \item Empirical evidence that residual predictive expertise and residual decision value diverge substantially on a real medical AI dataset, in a reward-dependent manner predicted by the theory.
  \item A synthetic allocation experiment showing that reward-aware review allocation outperforms uncertainty-based and expertise-based baselines under asymmetric utilities.
\end{enumerate}
